\chapter*{Введение}
\addcontentsline{toc}{chapter}{Введение}

\textbf{Актуальность исследования.} В последние годы автономная мобильная робототехника активно внедряется в различные сферы жизни человека, начиная от логистических складов и заканчивая обслуживанием в быту и ритейле. Одной из главных проблем автономных сервисных роботов является их функционирование в высокодинамичных, часто непредсказуемых средах, где объекты (люди, предметы мебели, товары) постоянно перемещаются или меняют свои состояния. Традиционные методы SLAM (Simultaneous Localization and Mapping) и 2D/3D картографирования отлично справляются с геометрической навигацией, но не обладают семантическим пониманием среды на уровне, необходимом для выполнения высокоуровневых команды, привязанных к контексту и времени (например, ``принеси мне чашку, которая пять минут назад была на столе в кухне'').

В связи с этим все больший интерес исследователей привлекают системы, объединяющие семантическое моделирование с помощью темпоральных графов знаний (Temporal Knowledge Graphs, TKG) и анализ естественного языка посредством больших языковых моделей (Large Language Models, LLM) \cite{ref7}. Темпоральные графы позволяют структурированно хранить многомерную информацию об объектах сцены, их свойствах и отношениях в различные моменты времени (исторический контекст). В то же время LLM дают возможность интерпретировать сложные составные команды человека и применять механизмы агентного рассуждения (Agentic Deep Graph Reasoning) для планирования действий робота \cite{ref1, ref19, ref20}.

Несмотря на активное развитие методов пополнения временных графов знаний (TKGC) и нейросимволического графового вывода, \textbf{на данный момент не существует готовых программных средств и фреймворков, позволяющих эффективно и инкрементально обновлять темпоральный граф знаний в режиме реального времени, которым мог бы пользоваться мобильный робот для навигации в изменяющейся среде.} Большинство существующих алгоритмов работы с графами знаний либо статичны, либо требуют полной перестройки структуры при внесении новых данных, либо не учитывают пространственно-временную специфику навигации робота. Поэтому разработка инкрементально обновляемой архитектуры графа знаний и программного модуля управления на её основе является актуальной научной и инженерной задачей.

\textbf{Цель работы:} Разработать программный модуль интеллектуального управления, обеспечивающий способность мобильного (сервисного) робота формировать, контролировать и адаптировать планы выполнения составных задач в условиях изменяющейся среды, используя инкрементально обновляемые темпоральные графы знаний для моделирования контекста и LLM для семантического анализа.

\textbf{Объект исследования:} Системы интеллектуального планирования действий и навигации сервисных роботов в динамических средах.

\textbf{Предмет исследования:} Инкрементальное обновление темпоральных графов знаний и методы агентного рассуждения (Agentic Reasoning) на базе больших языковых моделей для семантического контроля действий робота.
